\clearpage
\section*{September 30, 2026: What splitters look like, and starting lots}
\textit{Agent-drafted research entry.}

The cost of splitting is the least settled part of the model: specifications
that fit about as well give 761 to 1,697 units lost. Jacob asked what splitters
looked like before and after 485-x. The exploration used A/B rental parents of
at least 50 units, 1,173 first filed in 2014--2022 and 319 in 2025--2026.

\paragraph{Before 2025.} About 13 percent of parents split. Splitters were
larger (31 percent of parents of 300 or more units) and on larger, multi-lot
sites. Lot frontage and depth did not separate them, and an architect's
splitting on other projects barely predicted it. A third put several buildings
on one tax lot, 46 percent listed different owner names, and the buildings were
filed a median 11 days apart.

\paragraph{After.} Recent splits look like threshold avoidance: 91 percent put
each building on its own tax lot, 92 percent keep one owner, the buildings are
filed a median 3 days apart, a third follow a DOF lot subdivision near the first
filing (a fifth before 2025), and only 8 percent follow a lot merger (a quarter
before). The eligible site of 485-x is a tax lot, or a zoning lot with several
buildings in one application. Among parents of 100--299 units, 50--72 percent
split, against 9--15 percent before. Architects who never split before 2025
split as often as those who did.

\paragraph{Starting lots.} \path{build_parent_site_characteristics} now records
\texttt{starting\_lots}: the tax lots a site occupied before its own mergers and
subdivisions. From the lots of the parent's buildings at filing, each DOF merger
or subdivision from two years before the first filing to 180 days after it is
undone, in both periods; a reviewed parent keeps its documented parcels. One
and three years give nearly the same counts. The parent panels and the
estimation sample carry it; no other column changed. In the main sample, 26
percent of historical and 28 percent of recent parents start on two or more
lots.

\begin{center}
\small
\begin{tabular}{lrrrr}
\hline
Parents of 90+ units & One building at 99 & Split & One building 100--299 & Parents \\
\hline
Recent, one starting lot & 31\% & 10\% & 42\% & 106 \\
Recent, two or more & 6\% & 83\% & 9\% & 70 \\
Before 2025, one starting lot & 0\% & 9\% & 73\% & 489 \\
Before 2025, two or more & 0\% & 34\% & 53\% & 201 \\
\hline
\end{tabular}
\end{center}

On a single starting lot splitting hardly changed; those developers shrink to
99. On several lots they split and rarely shrink. The mean splitting cost
should therefore depend on starting lots, and units lost depend on how many
sites start on one lot.

\paragraph{Small parents.} Parents of 50--89 units split less after the policy
(4.6 percent against 7.6 percent in 2019--2022 on a common 180-day window,
$p = 0.47$), but the share had fallen since 2014--2018 (9.4 percent), and 2026
parents lack follow-up. Nothing in 485-x treats a small split differently: any
site of 6--99 units is a modest rental project with the same terms, and both
programs exempt buildings under 30 units from the building-service wage. Jacob
accepts that the model may not match this secular decline. Under 421-a(16),
rental projects of 300 or more units in the Enhanced Affordability Areas owed a
construction wage, which may explain the frequent splitting of large historical
parents.

\paragraph{Soltas (2024).} His study of the older 421-a program (2003--2015)
takes buildings as fixed: its incentive was roughly proportional to building
value, and he finds tax planning, not construction, the main response. A notch
at a unit count is what makes size respond here. The historical sample spans
that program (2014--2015), the 2016--2017 lapse, and Affordable New York
(2017--2022); the main 2019--2022 benchmark lies within the last. He also reads
421-a receipt from DOF exemption codes, which could check separate assessment
of historical splitters.

\paragraph{Follow-up: the splitting cost by starting lots.}
\path{fit_heterogeneity.R} now fits the main model with recent parents on one
starting lot and on several compared separately, each against the historical
parents of the same group, and with mean splitting cost
$\sigma e^{-\beta}$ on sites of two or more starting lots
(\path{scaled_burden_lot_count_splitting}). Fit by group, the main model
($\beta = 0$, \texttt{scaled\_burden\_by\_lots}) has log-likelihood $-570$; the
variant has $-553$ ($\beta = 4$), a likelihood ratio of 34 for one parameter.
Its mean splitting cost on one lot is 2.0, ten times the main model's, and
splitting on several lots is about 55 times cheaper. It predicts splitting by 6.9
percent of recent parents on one lot and 55.7 percent on several, against 5.5
and 55.1 observed; the main model predicts 17.7 and 46.5.

\paragraph{Why the estimate rises.} The pooled likelihood cannot tell shrinking
from splitting. The main estimate (761 units lost), the main model fit by group
(1,069) and the variant (1,717) lie within 3 log points of each other on the
pooled cells. The main estimate reaches 761 by predicting that single-lot sites
split three times as often as they do. Fit by group with one splitting cost,
the model compromises on a costlier cost everywhere ($\sigma$ 0.20 to 0.32,
$\gamma$ 2 to 3), only 2 log points better than the main estimate; the variant
improves on it by 17.

\paragraph{Bootstrap.} \path{bootstrap_scaled_burden.R} takes the model as a
second argument; the main model is its case of one group and $\beta = 0$. Over
500 draws of the main sample the variant gives:

\begin{center}
\small
\begin{tabular}{lrr}
\hline
 & Estimate & Bootstrap 95\% interval \\
\hline
$\beta$ & 4 & 2--5 \\
Mean splitting cost on one lot $\sigma$ & 2.0 & 0.63--7.9 \\
Median jump $\kappa$ & 0.30 & 0.05--1.0 \\
Units lost & 1,717 & 903--3,507 \\
\hline
\end{tabular}
\end{center}

No draw sets $\beta$ below 1.5, and 1 percent of draws give fewer units lost
than the main estimate; the bootstrap median is 1,843. Starting lots settle
that single-lot sites rarely split, but not how costly splitting them is or how
large the burden is, which together decide how many of them shrink. Units lost
is lowest near $\beta = 3$ (median 1,557 among those draws) and higher at 2 and
5. The variant is not yet the main model.

\paragraph{Follow-up, October 1: fragments and the recorded zoning lot.}
Many recent single-lot parents at 99 units sit on lots too small for their
units (99 units on 1,700--3,200 square feet). The filing lot is one piece of the
site: a zoning lot description recorded in ACRIS lists every tax lot of a
development's zoning lot, and for these parents it often lists several. Some
are lots built on, assembled but not yet merged; others are neighbours whose
buildings remain and lend floor area. Two new acquisition tasks publish the
records and the evidence for telling them apart:
\path{fetch_acris_zoning_lots} (31,943 zoning lot descriptions and 790
development-rights transfers since 2012, with the lots they list) and
\path{fetch_dob_demolition_filings} (BIS and DOB NOW full demolitions).
\path{build_parent_site_characteristics} adds, alongside the established land
measures, the site lots, land, residential FAR and frontage of the recorded
zoning lot: a listed lot is built on when vacant before filing or demolished
around it, and otherwise lends floor area. A record counts when it lists a lot
of the parent's buildings and is dated from two years before the first filing
to 180 days after it.

Counted this way, several lots no longer guarantee splitting: among recent
parents of 90 or more units on several site lots, 14 percent stay at 99 in one
building (6 percent by starting lots). Single-lot sites still rarely split.
Among parents with a recorded zoning lot, fit by site lots, $\beta = 2.5$
improves the fit over $\beta = 0$ by a likelihood ratio of 8.3 and raises units
lost from 368 to 569, less than starting lots raised them on the full sample.

\paragraph{A flaw in the 180-day rule.} Records exist for 61 percent of
2019--2022 parents in the window but 47 percent of recent ones, and recording
delays do not explain it: documents reach ACRIS within weeks, and recent
coverage does not rise with longer follow-up. Coverage moves from year to year
(45--80 percent by half-year) and peaks for parents first filed from July 2021
to June 2022, the rush before 421-a expired. Rush parents eventually have a
record about as often as others (89 against 81 percent) but recorded it twice
as fast (median 2.3 against 4.8 months), so 71 against 46 percent fall within
180 days. The rule therefore measures how fast the paperwork was done. Over 12
months, coverage is nearly equal for parents observed that long (69 against 65
percent). The window is to be redone before these measures enter the model;
the fits with the site lots on the full sample and without the rush were
stopped.
